\chapter{Functoriality of the Triple Identification in \(\mathsf{A}\) and \(g\)}
\label{v4-arith:w9-r5:chapter}

\emph{Chapter~\ref{v4-arith:w8-a:deninger-chiral-homology}
established the triple identification
\(H^{\bullet}_{\Delta}=H^{\bullet}_{\mathrm{ch}}=\partial\text{-chiral}\)
on the primary Heisenberg scope for the baseline genus \(g=0\) and the
fixed arithmetic Heisenberg \(\mathsf{H}^{\mathrm{Ar}}\). The
functoriality residuals of \Cref{v4-arith:w8-a:residual}(5,6) remain:
the three cohomology assignments extend to functors on a category
\(\mathsf{ArChAlg}\) of arithmetic chiral algebras on
\(\overline{\mathrm{Spec}\,\mathbb{Z}}\), and the triple identification
is a natural isomorphism of these functors, uniformly compatible with
the genus-\(g\) restriction and induction maps of
Chapter~\ref{v4-arith:all-genus-koszul}. The master theorem is
\Cref{v4-arith:w9-r5:thm:master-functoriality}; it rests on the
chain-level identification \(\Xi\) of
\Cref{v4-arith:w9-r5:thm:chain-identification}, the naturality diagram
\Cref{v4-arith:w9-r5:thm:naturality-A}, and the genus-stability theorem
\Cref{v4-arith:w9-r5:thm:genus-uniformity}. The diagonal morphism
\(\mathsf{H}^{\mathrm{Ar}}\to\mathsf{H}^{\mathrm{Ar}}\oplus\mathsf{H}^{\mathrm{Ar}}\)
is the explicit witness. All residuals of this chapter are inherited
from Chapters~\ref{v4-arith:3d-acs-E1bar} and
\ref{v4-arith:w8-a:deninger-chiral-homology}; no new residual is
introduced.}

\section{The Category of Arithmetic Chiral Algebras}
\label{v4-arith:w9-r5:category}

The three cohomology assignments of
Chapter~\ref{v4-arith:w8-a:deninger-chiral-homology} are defined on
an individual object, the arithmetic Heisenberg
\(\mathsf{H}^{\mathrm{Ar}}\). To state functoriality we introduce the
ambient category.

\begin{definition}[arithmetic chiral algebra]
\label{v4-arith:w9-r5:def:ArChAlg-object}
\ClaimStatusDefinitional.
An \emph{arithmetic chiral algebra on} \(\overline{\mathrm{Spec}\,\mathbb{Z}}\)
is a datum
\begin{equation}
\label{v4-arith:w9-r5:eq:ArChAlg-object}
\mathsf{A}
\;=\;
\Bigl(\mathsf{A}^{(\infty)},\,\{\mathsf{A}^{(p)}\}_{p\in\mathbb{P}},\,
\{u^{(v)}\}_{v\in|\overline{C}|},\,\{m^{(v,w)}\}_{v<w},\,
\{\partial^{\mathrm{PG}}_{v,w}\}_{v<w}\Bigr)
\end{equation}
subject to the axioms below. Write \(|\overline{C}|=\mathbb{P}\cup\{\infty\}\).
\begin{enumerate}
\item \emph{Archimedean factor.} \(\mathsf{A}^{(\infty)}\) is a
pointed \(E_{1}\)-chiral factorisation algebra on \(\mathbb{C}_{\infty}\)
in the Beilinson--Drinfeld / Costello--Gwilliam sense
(Ch.~\ref{v4-arith:platonic:def:S2}), together with a conformal vector
\(T\in\mathsf{A}^{(\infty)}_{2}\) and a spherical unit
\(u^{(\infty)}\in\mathsf{A}^{(\infty)}_{0}\).
\item \emph{Finite-prime factors.} For every prime \(p\in\mathbb{P}\),
\(\mathsf{A}^{(p)}\) is a pointed \(E_{0}\)-chiral coefficient object on
the Fargues--Scholze perfectoid diamond chart \(S_{p}\), equipped with
a Bost--Connes Adams endomorphism \(\psi_{p}\in\mathrm{End}(\mathsf{A}^{(p)})\)
and a Frobenius-semilinear orientation
\(\mathfrak{J}^{\mathrm{Ar}}_{p}\) (the finite-place substitute for the
archimedean \(i\in\mathbb{C}\)).
\item \emph{Spherical units.} The collection \(\{u^{(v)}\}_{v\in|\overline{C}|}\)
satisfies the Tate spherical compatibility
(\Cref{v4-arith:tate:setup:data}): \(u^{(v)}\) is closed for the local
bar differential \(d_{v}\), is a compact central idempotent, and is
the monoidal unit at almost every prime.
\item \emph{Adjacent products.} For every pair of distinct places
\(v<w\) in the auxiliary order \(2<3<5<\cdots<\infty\), the datum \(m^{(v,w)}\)
is a map
\[
m^{(v,w)}\colon\mathsf{A}^{(v)}\widehat{\otimes}\mathsf{A}^{(w)}
\longrightarrow\mathsf{A}^{(vw)},
\]
with \(\mathsf{A}^{(vw)}\) the formal merged-block coefficient in the
sense of \Cref{v4-arith:3d-acs:def:ordered-chiral-homology}. The
\(m^{(v,w)}\) are required to be \(\psi_{p}\)-equivariant at every finite
\(v=p\) and to intertwine the archimedean \(L_{0}\)-action at \(v=\infty\).
\item \emph{Parshin--Gersten boundary.} For every pair \(v<w\) of
ramified places, \(\partial^{\mathrm{PG}}_{v,w}\) is the Kato--Parshin
tame-symbol boundary
(\Cref{v4-arith:platonic:eq:bar-differential}); the collection
satisfies the Maurer--Cartan equation on every reciprocity-closed
finite stage
(\Cref{v4-arith:platonic:prop:finite-stage-strictification}).
\item \emph{Tate compatibility.} The restricted tensor product
\(\bigotimes'_{v}\mathsf{A}^{(v)}\) in pointed \(E_{2}\)-monoidal stable
presentable \(\infty\)-categories
(\Cref{v4-arith:tate:cor:F2-closed}) is well-defined, i.e.\ almost
every \(u^{(v)}\) is the monoidal unit and the Tate-pair
compatibility F2.c of \Cref{v4-arith:cl:F2-closure} is satisfied.
\end{enumerate}
The archimedean \(\mathsf{A}^{(\infty)}\) is always genuinely chiral;
finite factors may be \(E_{0}\) per
\Cref{v4-arith:3d-acs:def:ordered-chiral-homology}. The object
\(\mathsf{A}\) is \emph{abelian Heisenberg-type} if every
\(\mathsf{A}^{(p)}\) is a rank-\(2g\) bosonic Fock for some fixed \(g\geq 0\)
(the genus of \(\mathsf{A}\)), and \(\mathsf{A}^{(\infty)}\) is the
rank-\(2g\) archimedean Heisenberg of
\Cref{v4-arith:thm:heisenberg-bar-zeta}.
\end{definition}

\begin{remark}[on the merged-block target]
\label{v4-arith:w9-r5:rem:merged-block}
The formal merged-block coefficient \(\mathsf{A}^{(vw)}\) is not a new
arithmetic place. It is the target of the adjacent contraction map
in the sense of the ordered incidence category
\(\mathsf{Inc}_{<}(V)\) of
\Cref{v4-arith:3d-acs:def:ordered-incidence-category}; concretely a
completed tensor product of \(\mathsf{A}^{(v)}\) and \(\mathsf{A}^{(w)}\)
modulo a quotient relation recording that the two local operators
act on the same global fibre. This convention is essential to avoid
the naive reading that primes are fused into composite arithmetic
objects.
\end{remark}

\begin{definition}[morphism of arithmetic chiral algebras]
\label{v4-arith:w9-r5:def:ArChAlg-morphism}
\ClaimStatusDefinitional.
A \emph{morphism} \(\phi\colon\mathsf{A}_{1}\to\mathsf{A}_{2}\) of
arithmetic chiral algebras is a collection
\(\phi=\{\phi^{(v)}\}_{v\in|\overline{C}|}\) with
\(\phi^{(v)}\colon\mathsf{A}_{1}^{(v)}\to\mathsf{A}_{2}^{(v)}\) such
that
\begin{enumerate}
\item \(\phi^{(\infty)}\) is an \(E_{1}\)-chiral factorisation map
preserving the conformal vector;
\item at every prime \(p\), \(\phi^{(p)}\) is \(\psi_{p}\)-equivariant and
\(\mathfrak{J}^{\mathrm{Ar}}_{p}\)-compatible: \(\phi^{(p)}\circ\psi_{p,1}=\psi_{p,2}\circ\phi^{(p)}\);
\item the spherical units are preserved: \(\phi^{(v)}(u_{1}^{(v)})=u_{2}^{(v)}\)
for every \(v\in|\overline{C}|\);
\item the adjacent products are intertwined:
\(m_{2}^{(v,w)}\circ(\phi^{(v)}\otimes\phi^{(w)})=\phi^{(vw)}\circ m_{1}^{(v,w)}\)
for every \(v<w\);
\item the Parshin--Gersten boundaries are intertwined:
\(\phi\circ\partial_{1}^{\mathrm{PG}}=\partial_{2}^{\mathrm{PG}}\circ\phi\)
at every pair of ramified places;
\item the restricted-tensor map
\(\widehat{\phi}:=\bigotimes'_{v}\phi^{(v)}\) is a map in the pointed
\(E_{2}\)-monoidal category of
\Cref{v4-arith:tate:cor:F2-closed}; this is automatic from (1)--(5)
at almost every place by the spherical-unit preservation.
\end{enumerate}
Composition is place-wise; identities are place-wise identities. The
resulting category is \(\mathsf{ArChAlg}\).
\end{definition}

\begin{proposition}[\(\mathsf{ArChAlg}\) is a well-defined category]
\label{v4-arith:w9-r5:prop:ArChAlg-is-category}
\ClaimStatusProvedHere. Composition of morphisms satisfies the
axioms of \Cref{v4-arith:w9-r5:def:ArChAlg-morphism}: axioms (1)--(5)
are place-wise, and axiom (6) follows from place-wise axioms by the
almost-every-place triviality of the units. Identity morphisms are
identities in every axiom. Associativity and unitality are
inherited from the underlying categories of \(E_{1}\)-chiral
factorisation algebras at \(\infty\) and \(E_{0}\)-coefficient categories
at finite places.
\end{proposition}

\begin{proof}
For composition: given \(\phi\colon\mathsf{A}_{1}\to\mathsf{A}_{2}\)
and \(\chi\colon\mathsf{A}_{2}\to\mathsf{A}_{3}\), the collection
\((\chi\circ\phi)^{(v)}:=\chi^{(v)}\circ\phi^{(v)}\) preserves each
axiom: axiom (1) is closed under composition of \(E_{1}\)-chiral maps;
axiom (2) is closed because \(\psi_{p}\)-equivariance composes; axiom
(3) uses \(\chi^{(v)}(\phi^{(v)}(u_{1}^{(v)}))=\chi^{(v)}(u_{2}^{(v)})=u_{3}^{(v)}\);
axiom (4) similarly using the stacked naturality diagrams. Axiom
(5) uses the linear action of the boundary. Associativity and
unitality are automatic at each place. Axiom (6) follows from the
almost-every-place compatibility of the Tate restricted tensor
(\Cref{v4-arith:cl:F2-closure}): the set of places at which
\(u^{(v)}\) fails to be a strict unit is bounded, and composition
preserves this bound.
\end{proof}

\begin{definition}[genus filtration]
\label{v4-arith:w9-r5:def:genus-filtration}
\ClaimStatusDefinitional. For every \(g\geq 0\), let
\(\mathsf{ArChAlg}_{g}\subset\mathsf{ArChAlg}\) be the full subcategory
of abelian Heisenberg-type arithmetic chiral algebras of genus \(g\)
in the sense of \Cref{v4-arith:w9-r5:def:ArChAlg-object}. The
canonical inclusion
\(\iota_{g,g+1}\colon\mathsf{ArChAlg}_{g}\hookrightarrow\mathsf{ArChAlg}_{g+1}\)
is by tensoring with the rank-\(2\) archimedean Heisenberg
(\Cref{v4-arith:ag:defn:datum}'s rank-\(2\) enhancement at each
place). The retraction
\(\pi_{g+1,g}\colon\mathsf{ArChAlg}_{g+1}\to\mathsf{ArChAlg}_{g}\)
projects onto the first \(2g\) oscillators.
\end{definition}

\begin{remark}[why the filtration]
\label{v4-arith:w9-r5:rem:filtration-reason}
The filtration records that the genus-\(g\) category embeds into the
genus-\((g+1)\) category in a canonical way (tensor with a fresh
rank-\(2\) factor at every place), and that one can restrict to any
subgenus. The adjunction \((\iota,\pi)\) is a \emph{reflective
inclusion}: \(\pi_{g+1,g}\circ\iota_{g,g+1}=\mathrm{id}_{g}\) is strict,
and the counit \(\iota_{g,g+1}\circ\pi_{g+1,g}\Rightarrow\mathrm{id}_{g+1}\)
is the natural transformation that kills the freshly-added rank-\(2\)
factor. The pair \((\iota_{g,g+1},\pi_{g+1,g})\) is the canonical reflective
pair through which all genus functoriality statements below pass.
\end{remark}

\begin{remark}[source anchors and disjointness]
\label{v4-arith:w9-r5:rem:source-anchors}
The archimedean \(E_{1}\)-chiral structure is sourced from Beilinson--Drinfeld
2004 and Francis--Gaitsgory 2012 (see
\Cref{v4-arith:app:A:beilinson-etc}); the finite-place \(E_{0}\)-coefficient
structure is sourced from Fargues--Scholze arXiv:2102.13459 and
Bhatt--Scholze 2022 (prismatic; see \Cref{v4-arith:app:A:prismatic}).
These two source families are disjoint in the Vol~IV
sense. The morphism axioms (1)--(6) each arise from exactly one of
these source families; this disjointness is what makes the
functoriality statements below decorable in the sense of the
Vol~IV HZ-IV campaign.
\end{remark}

\section{The Three Functors}
\label{v4-arith:w9-r5:three-functors}

The cohomology assignments of
Chapter~\ref{v4-arith:w8-a:deninger-chiral-homology} are defined on the
single object \(\mathsf{H}^{\mathrm{Ar}}\). Functoriality requires a
target category. Write \(\mathrm{grVect}_{\mathbb{C}}^{\mathrm{Frob}}\)
for the category of finite-dimensional graded \(\mathbb{C}\)-vector spaces
equipped with a graded-linear automorphism \(F\) (the Frobenius) and a
graded-symmetric non-degenerate pairing \(\mathrm{PD}\) (the Poincar\'e
duality); morphisms are graded linear maps commuting with both structures.

\begin{definition}[Deninger functor \(H^{\bullet}_{\Delta}\)]
\label{v4-arith:w9-r5:def:functor-Delta}
\ClaimStatusDefinitional.
The \emph{Deninger functor}
\begin{equation}
\label{v4-arith:w9-r5:eq:functor-Delta}
H^{\bullet}_{\Delta}\colon
\mathsf{ArChAlg}\longrightarrow
\mathrm{grVect}_{\mathbb{C}}^{\mathrm{Frob}}
\end{equation}
sends an arithmetic chiral algebra \(\mathsf{A}\) to the triple
\((H^{\bullet}_{\Delta}(\mathsf{A}),F_{\mathsf{A}},\mathrm{PD}_{\mathsf{A}})\)
of \Cref{v4-arith:w8-a:def:deninger-cohomology}, with
\(F_{\mathsf{A}}=F_{\infty,\mathsf{A}}\otimes\bigotimes'_{p}\psi_{p,\mathsf{A}}\)
the global Frobenius of \Cref{v4-arith:w8-a:def:F-global} and
\(\mathrm{PD}_{\mathsf{A}}\) the Poincar\'e duality pairing of
\Cref{v4-arith:w8-a:thm:triple-duality}. On a morphism
\(\phi\colon\mathsf{A}_{1}\to\mathsf{A}_{2}\), the functor sends
\(\phi\) to the induced map
\(H^{\bullet}_{\Delta}(\phi)\colon H^{\bullet}_{\Delta}(\mathsf{A}_{1})\to H^{\bullet}_{\Delta}(\mathsf{A}_{2})\).
\end{definition}

\begin{definition}[chiral-homology functor \(H^{\bullet}_{\mathrm{ch}}\)]
\label{v4-arith:w9-r5:def:functor-ch}
\ClaimStatusDefinitional.
The \emph{chiral-homology functor}
\begin{equation}
\label{v4-arith:w9-r5:eq:functor-ch}
H^{\bullet}_{\mathrm{ch}}\colon
\mathsf{ArChAlg}\longrightarrow
\mathrm{grVect}_{\mathbb{C}}^{\mathrm{Frob}}
\end{equation}
sends \(\mathsf{A}\) to
\(H^{\bullet}(B^{\mathrm{Ar}}(\mathsf{A}),D_{\mathsf{A}})\) of
\Cref{v4-arith:platonic:def:bar}, with \(D_{\mathsf{A}}\) the
signed total differential determined by \(\mathsf{A}\). On a morphism
\(\phi\colon\mathsf{A}_{1}\to\mathsf{A}_{2}\), the induced chain map
\(B^{\mathrm{Ar}}(\phi)\) commutes with \(D_{\mathsf{A}_{1}}\) and
\(D_{\mathsf{A}_{2}}\) by axioms (3), (4), (5) of
\Cref{v4-arith:w9-r5:def:ArChAlg-morphism}; pass to cohomology for
\(H^{\bullet}_{\mathrm{ch}}(\phi)\).
\end{definition}

\begin{definition}[bulk-boundary functor pair \(\partial\circ\mathrm{BdryACS}\)]
\label{v4-arith:w9-r5:def:functor-bdryACS}
\ClaimStatusDefinitional.
The \emph{bulk-boundary assignment} is the composite
\begin{equation}
\label{v4-arith:w9-r5:eq:functor-bdryACS}
\partial\circ\mathrm{BdryACS}\colon
\mathsf{ArChAlg}
\xrightarrow{\mathrm{BdryACS}}
\mathsf{3DACS}_{\overline{\mathrm{Spec}\,\mathbb{Z}}}
\xrightarrow{\partial}
\mathrm{grVect}_{\mathbb{C}}^{\mathrm{Frob}},
\end{equation}
where
\begin{itemize}
\item \(\mathsf{3DACS}_{\overline{\mathrm{Spec}\,\mathbb{Z}}}\) is the
category of conjectural 3D quantum arithmetic Chern--Simons theories
on the \'etale 3-space \(\overline{\mathrm{Spec}\,\mathbb{Z}}\) in the
sense of \Cref{v4-arith:3d-acs:conj:qu-acs}, with morphisms given
by bulk interface defects;
\item \(\mathrm{BdryACS}\) is the canonical lift that takes an arithmetic
chiral algebra to the 3D bulk theory whose boundary WZW is
\(\mathsf{A}^{(\infty)}\) and whose finite-place defects are labelled
by \(\{\mathsf{A}^{(p)}\}\) via the Morishita knot-prime dictionary
(\Cref{v4-arith:3d-acs:cor:finite-wilson},
\Cref{v4-arith:w8-a:lem:wilson-finite});
\item \(\partial\) is the boundary chiral homology functor, sending a
3D bulk theory to the chiral homology of its boundary factorisation
algebra on the archimedean fibre.
\end{itemize}
On a morphism \(\phi\colon\mathsf{A}_{1}\to\mathsf{A}_{2}\),
\(\mathrm{BdryACS}\) produces a bulk interface defect between the two
3D theories whose boundary restriction is \(\phi^{(\infty)}\) and
whose finite-place Wilson-line restriction is \(\{\phi^{(p)}\}\).
\end{definition}

\begin{remark}[conditional status of the bulk-boundary lift]
\label{v4-arith:w9-r5:rem:bdryACS-conditional}
The functor \(\mathrm{BdryACS}\) is conditional on the 3D quantum ACS
existence conjecture \Cref{v4-arith:3d-acs:conj:qu-acs}. The boundary
functor \(\partial\) is conditional on the local-trace comparison
\Cref{v4-arith:3d-acs:thm:bulk-boundary-Vprim}. Both conditionals are
inherited from Chapter~\ref{v4-arith:3d-acs-E1bar} and are explicitly
carried through each theorem of this chapter.
\end{remark}

\begin{remark}[source disjointness of the three functors]
\label{v4-arith:w9-r5:rem:three-functors-disjoint}
The Deninger functor is derived from Deninger 1998 (arithmetic
absolute-value axioms). The chiral-homology functor is derived from
the Platonic bar complex of Beilinson--Drinfeld + Francis
(\Cref{v4-arith:platonic:def:bar}). The bulk-boundary functor is
derived from Kim 2015 arXiv:1510.05818 and the conjectural quantum
ACS of \Cref{v4-arith:3d-acs:conj:qu-acs}. These three derivation
paths are pairwise source-disjoint; the natural isomorphism theorem
below is a genuine independent verification of the triple equality,
not a restatement of any single derivation.
\end{remark}

\section{Chain-Level Identification}
\label{v4-arith:w9-r5:chain-level}

The three functors land in \(\mathrm{grVect}_{\mathbb{C}}^{\mathrm{Frob}}\)
by taking cohomology of chain complexes that are not obviously the same.
The natural isomorphism follows once one identifies the chain complexes
before passing to cohomology.

\begin{definition}[universal arithmetic bar]
\label{v4-arith:w9-r5:def:universal-bar}
\ClaimStatusDefinitional. For \(\mathsf{A}\in\mathsf{ArChAlg}\), the
\emph{universal arithmetic bar} of \(\mathsf{A}\) is the pair
\begin{equation}
\label{v4-arith:w9-r5:eq:universal-bar}
\mathcal{B}^{\mathrm{univ}}(\mathsf{A})
\;:=\;\bigl(B^{\mathrm{ord}}_{\bullet}(\mathsf{A}),
B^{\mathrm{Ar}}_{\bullet}(\mathsf{A})\bigr),
\end{equation}
with \(B^{\mathrm{ord}}_{\bullet}\) the ordered incidence bar of
\Cref{v4-arith:3d-acs:defn:E1-bar} applied to the collection
\(\{\mathsf{A}^{(v)}\}_{v\in|\overline{C}|}\), and \(B^{\mathrm{Ar}}_{\bullet}\)
the Platonic arithmetic-chiral bar of
\Cref{v4-arith:platonic:def:bar}. Both are equipped with their
natural signed differentials.
\end{definition}

\begin{theorem}[chain-level identification on \(\mathsf{ArChAlg}\)]
\label{v4-arith:w9-r5:thm:chain-identification}
\ClaimStatusProvedHereConditional \emph{(on
\Cref{v4-arith:3d-acs:thm:chi-hom-bar} and
\Cref{v4-arith:platonic:prop:bar-comparison})}. For every
\(\mathsf{A}\in\mathsf{ArChAlg}\) there is a canonical quasi-isomorphism
\begin{equation}
\label{v4-arith:w9-r5:eq:chain-identification}
\Xi_{\mathsf{A}}\colon
B^{\mathrm{ord}}_{\bullet}(\mathsf{A})
\xrightarrow{\;\sim\;}
B^{\mathrm{Ar}}_{\bullet}(\mathsf{A})
\end{equation}
of chain complexes, and
\(\Xi\colon B^{\mathrm{ord}}_{\bullet}\Rightarrow B^{\mathrm{Ar}}_{\bullet}\)
is a natural transformation of functors
\(\mathsf{ArChAlg}\to\mathrm{Ch}_{\mathbb{C}}\). The induced maps on
cohomology coincide with the chiral-homology functor
\Cref{v4-arith:w9-r5:def:functor-ch}.
\end{theorem}

\begin{proof}
Construction of \(\Xi_{\mathsf{A}}\). Both complexes are defined as
filtered colimits over reciprocity-closed finite stages
\(S\in\mathrm{Fin}_{\mathrm{rec}}(|\overline{C}|)\). On a finite stage
\(S\) with an ordered block decomposition
\(B_{1}|\cdots|B_{r}\) of \(S\), the ordered incidence bar assigns the tensor product
\(\mathcal{F}_{\mathsf{A}}(B_{1}|\cdots|B_{r}):=\bigotimes_{i}\mathsf{A}^{(B_{i})}\)
(using \(\mathcal{F}\) to avoid collision with the Frobenius \(F_{\mathsf{A}}\))
with contractions given by the adjacent products \(m^{(v,w)}\); the
Platonic bar assigns
\(B^{\mathrm{ch}}(\mathsf{A}^{(\infty)})\widehat{\otimes}
C^{\bullet}_{\mathrm{PG},S}(\overline{C})\widehat{\otimes}
\widehat{\bigotimes}_{p\in S_{f}}B^{\Lambda_{p}}_{E_{1}}(\mathsf{A}^{(p)})\).

At the archimedean place, the \(E_{1}\)-chiral bar of
\(\mathsf{A}^{(\infty)}\) is the \(N_{\bullet}\)-normalisation of the
semi-simplicial representation (\(\operatorname{srep}\): the functor
from the semi-simplicial nerve to chain complexes via the
Dold--Kan correspondence) of the
\(E_{1}\)-nerve restricted to singleton blocks; this is the content of
\Cref{v4-arith:platonic:prop:bar-comparison} on the archimedean
factor alone. Call this isomorphism \(\Xi^{(\infty)}\).

At a finite prime \(p\), the ordered-incidence block contribution is
\(\mathsf{A}^{(p)}\) itself (singleton block, no internal subdivision
because \(\mathsf{A}^{(p)}\) is \(E_{0}\)); the Platonic finite factor
is \(B^{\Lambda_{p}}_{E_{1}}(\mathsf{A}^{(p)})\), the \(\Lambda_{p}\)-equivariant
bar of the local Heisenberg. These are quasi-isomorphic by the
arithmetic Positselski bar-cobar adjunction
(\Cref{v4-arith:thm:arithmetic-positselski}) composed with the
triviality of the \(E_{1}\)-bar of an \(E_{0}\)-algebra: the bar of a
unital \(E_{0}\)-object reduces to the shifted augmentation complex,
which is homologically the identity. Call this isomorphism
\(\Xi^{(p)}\).

The Parshin--Gersten layer \(C^{\bullet}_{\mathrm{PG},S}(\overline{C})\)
contributes at the ordered-incidence side through the merger maps
\(m^{(v,w)}\) with their tame-symbol boundaries; the match is
\Cref{v4-arith:platonic:prop:finite-stage-strictification}. Put
\(\Xi_{S}:=\Xi^{(\infty)}\otimes(\bigotimes_{p\in S_{f}}\Xi^{(p)})\otimes\mathrm{id}_{\mathrm{PG},S}\);
the Maurer--Cartan closure at finite stage gives that \(\Xi_{S}\)
commutes with the two signed total differentials.

Passage to the colimit. The transition maps at both sides insert
closed spherical units; by \Cref{v4-arith:platonic:lem:closed-units}
the spherical units are killed by the local bar differentials on
both sides, so \(\Xi_{S}\) extends compatibly. The filtered colimit
\(\Xi:=\operatorname{colim}_{S}\Xi_{S}\) is the required
quasi-isomorphism.

Naturality in \(\mathsf{A}\). A morphism
\(\phi\colon\mathsf{A}_{1}\to\mathsf{A}_{2}\) induces a map of both
ordered-bar and Platonic-bar assignments (axioms (3), (4), (5) of
\Cref{v4-arith:w9-r5:def:ArChAlg-morphism}). The archimedean factor
map \(\Xi^{(\infty)}\) is natural because \(E_{1}\)-chiral bar is
functorial in \(\mathsf{A}^{(\infty)}\) (standard BD/CG). The
finite-prime factor map \(\Xi^{(p)}\) is natural because arithmetic
Positselski is functorial in the underlying \(\Lambda_{p}\)-module.
Both naturalities glue over the filtered stage system, giving that
\(\Xi\) is a natural transformation of functors
\(\mathsf{ArChAlg}\to\mathrm{Ch}_{\mathbb{C}}\).

Descent to cohomology. Applying \(H^{\bullet}\) to \(\Xi\) yields the
identity between the ordered-bar homology
(\Cref{v4-arith:3d-acs:eq:chi-hom-bar}) and the Platonic-bar
homology
(\Cref{v4-arith:w8-a:eq:chiral-homology-def}); both are the
chiral-homology functor by \Cref{v4-arith:w9-r5:def:functor-ch}.
\end{proof}

\begin{corollary}[factoring the three functors]
\label{v4-arith:w9-r5:cor:factoring}
\ClaimStatusProvedHereConditional \emph{(on
\Cref{v4-arith:w9-r5:thm:chain-identification})}. The three
functors of
\Cref{v4-arith:w9-r5:three-functors} all factor through the
chain-level homology functor
\(\mathsf{A}\mapsto H^{\bullet}(B^{\mathrm{Ar}}(\mathsf{A}),D_{\mathsf{A}})\):
\[
H^{\bullet}_{\Delta}(\mathsf{A})
\;\underset{\mathrm{def}}{=}\;H^{\bullet}_{\mathrm{ch}}(\mathsf{A})
\;=\;\partial\circ\mathrm{BdryACS}(\mathsf{A})
\]
on every \(\mathsf{A}\in\mathsf{ArChAlg}\), modulo the conditional
pieces of Chapters~\ref{v4-arith:3d-acs-E1bar} and
\ref{v4-arith:w8-a:deninger-chiral-homology}.
\end{corollary}

\begin{proof}
The left-hand equality is
\Cref{v4-arith:w8-a:def:deninger-cohomology} (definitional). The
right-hand equality follows from
\Cref{v4-arith:3d-acs:thm:chi-hom-bar} and
\Cref{v4-arith:w9-r5:thm:chain-identification}:
\(\partial\circ\mathrm{BdryACS}(\mathsf{A})\) computes the boundary
chiral homology of the 3D bulk theory attached to \(\mathsf{A}\),
which equals the ordered incidence homology of \(\{\mathsf{A}^{(v)}\}\),
which is quasi-isomorphic via \(\Xi_{\mathsf{A}}\) to the Platonic bar
homology \(H^{\bullet}_{\mathrm{ch}}(\mathsf{A})\).
\end{proof}

\subsection*{Source-disjointness and scope of the chain-level identification}
\label{v4-arith:w9-r5:ah:chain}

\emph{Source-disjointness.} The ordered-incidence bar rests on
Francis 2013 (geometric \(E_{1}\)-nerve); the Platonic bar rests on
Beilinson--Drinfeld 2004 and arithmetic Positselski
(\Cref{v4-arith:platonic:def:bar}). The two constructions come from
disjoint source families; \(\Xi\) is therefore a genuine
Vol~IV verification of the identification, not a tautology.

\emph{Signs.} The ordered-incidence differential uses alternating
contraction signs on simplicial nerves; the Platonic differential
\eqref{v4-arith:platonic:eq:bar-differential} uses Koszul signs on
tensor factors. These match under \(\Xi\) because the simplicial
normalisation formula turns alternating contraction signs into
Koszul tensor signs (standard Dold--Kan).

\emph{Ambient category.} Both bars take values in chain complexes
over \(\mathbb{C}\) with restricted-tensor refinement to the adelic
factorisation structure; the ambient category is identical.

\emph{Conditional hypothesis.} The quasi-isomorphism \(\Xi^{(p)}\)
on finite-prime factors uses arithmetic Positselski
\Cref{v4-arith:thm:arithmetic-positselski}, conditional on the
non-abelian Iwasawa extension. This is the only inherited
conditional; all other steps are unconditional within the
primary Heisenberg scope.

\emph{Naturality.} Axioms (3), (4), (5) of
\Cref{v4-arith:w9-r5:def:ArChAlg-morphism} supply precisely the
intertwining needed at each layer (units, adjacent products,
Parshin--Gersten boundary); no additional functoriality input is
required.

\emph{Unconditional scope.} The equivalence between ordered-bar
and Platonic-bar homology on the primary Heisenberg scope is the
content of \Cref{v4-arith:platonic:prop:bar-comparison}, which
holds unconditionally within that scope. No compute engine backs
this identification.

\section{Functoriality in \(\mathsf{A}\)}
\label{v4-arith:w9-r5:functoriality-A}

\Cref{v4-arith:w9-r5:cor:factoring} equates the three functors on
objects. A natural isomorphism requires that every morphism
\(\phi\colon\mathsf{A}_{1}\to\mathsf{A}_{2}\) in \(\mathsf{ArChAlg}\)
induces commuting diagrams; the commutation is the content of the
next theorem.

\begin{theorem}[naturality in \(\mathsf{A}\)]
\label{v4-arith:w9-r5:thm:naturality-A}
\ClaimStatusProvedHereConditional \emph{(on
\Cref{v4-arith:w9-r5:thm:chain-identification},
\Cref{v4-arith:3d-acs:conj:qu-acs}, and
\Cref{v4-arith:3d-acs:thm:bulk-boundary-Vprim})}. The three natural
isomorphisms
\begin{equation}
\label{v4-arith:w9-r5:eq:natural-iso}
H^{\bullet}_{\Delta}\;\xRightarrow{\;\eta_{\Delta,\mathrm{ch}}\;}\;H^{\bullet}_{\mathrm{ch}}
\;\xRightarrow{\;\eta_{\mathrm{ch},\partial}\;}\;\partial\circ\mathrm{BdryACS}
\end{equation}
are natural: for every morphism \(\phi\colon\mathsf{A}_{1}\to\mathsf{A}_{2}\)
of arithmetic chiral algebras, the three squares
\[
\begin{tikzcd}[column sep=large,row sep=large]
H^{\bullet}_{\Delta}(\mathsf{A}_{1})\arrow[r,"H^{\bullet}_{\Delta}(\phi)"]\arrow[d,"\eta_{\Delta,\mathrm{ch},\mathsf{A}_{1}}"']
& H^{\bullet}_{\Delta}(\mathsf{A}_{2})\arrow[d,"\eta_{\Delta,\mathrm{ch},\mathsf{A}_{2}}"]\\
H^{\bullet}_{\mathrm{ch}}(\mathsf{A}_{1})\arrow[r,"H^{\bullet}_{\mathrm{ch}}(\phi)"]\arrow[d,"\eta_{\mathrm{ch},\partial,\mathsf{A}_{1}}"']
& H^{\bullet}_{\mathrm{ch}}(\mathsf{A}_{2})\arrow[d,"\eta_{\mathrm{ch},\partial,\mathsf{A}_{2}}"]\\
\partial\circ\mathrm{BdryACS}(\mathsf{A}_{1})\arrow[r,"\partial\circ\mathrm{BdryACS}(\phi)"']
& \partial\circ\mathrm{BdryACS}(\mathsf{A}_{2})
\end{tikzcd}
\]
commute. The three horizontal maps are identified under the
naturality, and the induced Frobenius and Poincar\'e duality
structures are preserved pairwise: \(F_{\mathsf{A}_{2}}\circ H^{\bullet}(\phi)=H^{\bullet}(\phi)\circ F_{\mathsf{A}_{1}}\)
and likewise for \(\mathrm{PD}\).
\end{theorem}

\begin{proof}
The upper square. The equality
\(H^{\bullet}_{\Delta}(\mathsf{A})\underset{\mathrm{def}}{=}H^{\bullet}_{\mathrm{ch}}(\mathsf{A})\)
is definitional (\Cref{v4-arith:w8-a:def:deninger-cohomology}); the
natural transformation \(\eta_{\Delta,\mathrm{ch}}\) is the identity.
Naturality is immediate: the identity intertwines every induced map.

The lower square. Apply \Cref{v4-arith:w9-r5:thm:chain-identification}
to both \(\mathsf{A}_{1}\) and \(\mathsf{A}_{2}\), obtaining chain-level
quasi-isomorphisms \(\Xi_{\mathsf{A}_{1}}\) and \(\Xi_{\mathsf{A}_{2}}\).
The natural transformation \(\eta_{\mathrm{ch},\partial}\) is the
induced map on cohomology:
\(\eta_{\mathrm{ch},\partial,\mathsf{A}}:=\sigma_{\mathsf{A}}\circ H^{\bullet}(\Xi_{\mathsf{A}})^{-1}\),
where \(\sigma_{\mathsf{A}}\colon H^{\bullet}(B^{\mathrm{ord}}_{\bullet}(\mathsf{A}))\xrightarrow{\sim}\partial\circ\mathrm{BdryACS}(\mathsf{A})\)
is the isomorphism of \Cref{v4-arith:3d-acs:thm:chi-hom-bar}; this
composite is well-defined by that theorem and
\Cref{v4-arith:w9-r5:thm:chain-identification}.

Commutativity at the chain level. The induced chain map
\(B^{\mathrm{ord}}_{\bullet}(\phi)\) intertwines \(\Xi_{\mathsf{A}_{1}}\)
and \(\Xi_{\mathsf{A}_{2}}\) by the naturality of
\Cref{v4-arith:w9-r5:thm:chain-identification} in \(\mathsf{A}\). The
induced chain map \(B^{\mathrm{Ar}}(\phi)\) likewise intertwines, by
axioms (3), (4), (5) of \Cref{v4-arith:w9-r5:def:ArChAlg-morphism}:
the spherical units are mapped to spherical units; the adjacent
products commute through \(\phi\); the Parshin--Gersten boundary
commutes through \(\phi\). Composing, we get
\(\Xi_{\mathsf{A}_{2}}\circ B^{\mathrm{ord}}(\phi)=B^{\mathrm{Ar}}(\phi)\circ\Xi_{\mathsf{A}_{1}}\)
as maps of chain complexes. Passing to cohomology gives the
commutative square.

Frobenius compatibility. The global Frobenius
\(F_{\mathsf{A}}=F_{\infty,\mathsf{A}}\otimes\bigotimes'_{p}\psi_{p,\mathsf{A}}\)
commutes with \(H^{\bullet}(\phi)\) because axiom (2) of the morphism
definition supplies \(\psi_{p}\)-equivariance at each finite prime,
and axiom (1) supplies preservation of the conformal vector at
\(\infty\) (from which \(F_{\infty}\) is conjugated via
\Cref{v4-arith:w8-a:lem:L0-conjugacy}). Hence
\(F_{\mathsf{A}_{2}}\circ H^{\bullet}(\phi)=H^{\bullet}(\phi)\circ F_{\mathsf{A}_{1}}\).

Poincar\'e duality compatibility. The PD pairing is determined by
the Koszul self-duality identity
\Cref{v4-arith:w8-a:thm:triple-duality}. A morphism \(\phi\) in
\(\mathsf{ArChAlg}\) preserves the spherical units and the adjacent
products; the self-duality pairing is expressed chain-level as the
Positselski contraction pairing, which is natural in the underlying
\(\mathsf{A}\). Hence \(\mathrm{PD}\) commutes with \(\phi\).
\end{proof}

\begin{corollary}[triple functoriality in \(\mathsf{A}\)]
\label{v4-arith:w9-r5:cor:triple-A}
\ClaimStatusProvedHereConditional \emph{(on
\Cref{v4-arith:w9-r5:thm:naturality-A})}. The three functors
\(H^{\bullet}_{\Delta}\), \(H^{\bullet}_{\mathrm{ch}}\), and
\(\partial\circ\mathrm{BdryACS}\) are naturally isomorphic as
functors \(\mathsf{ArChAlg}\to\mathrm{grVect}_{\mathbb{C}}^{\mathrm{Frob}}\).
Both structural data (Frobenius, Poincar\'e duality) pull through
the natural isomorphisms.
\end{corollary}

\begin{proof}
Immediate from \Cref{v4-arith:w9-r5:thm:naturality-A} after recording
that \(\eta_{\Delta,\mathrm{ch}}\) is invertible (as the identity)
and \(\eta_{\mathrm{ch},\partial}\) is invertible as the induced
map on cohomology from a chain-level quasi-isomorphism.
\end{proof}

\subsection*{Source-disjointness and scope of naturality in \(\mathsf{A}\)}
\label{v4-arith:w9-r5:ah:A}

\emph{Source-disjointness.} Beilinson--Drinfeld 2004 supplies the
\(E_{1}\)-chiral functoriality at the archimedean factor; Kim 2015 and
the quantum ACS supply the bulk interface defect functoriality at the
3D side. These derivation paths are disjoint; the natural isomorphism
\(\eta_{\mathrm{ch},\partial}\) verified against both is a genuine
independent check.

\emph{Signs.} The bar-differential Koszul signs agree on both sides by
construction of \Cref{v4-arith:platonic:def:bar}; the
ordered-incidence alternating signs match under \(\Xi\) by the
simplicial normalisation formula.

\emph{Ambient category.} Both sides land in graded \(\mathbb{C}\)-vector
spaces with Frobenius and Poincar\'e duality; no compatibility issue.

\emph{Frobenius hypothesis.} The naturality of \(F_{\infty}\) under
\(\phi^{(\infty)}\) requires that \(\phi^{(\infty)}\) preserves the
conformal vector; this is axiom~(1) of
\Cref{v4-arith:w9-r5:def:ArChAlg-morphism}, which is part of the
definition of a morphism and imposes no hidden constraint.

\emph{Axiom independence.} Axioms (2) and (4) are logically
independent: axiom~(4) applied at \(v=p, w=\infty\) implies axiom~(2)
only on the unit line (the spherical vacuum at \(\infty\) is fixed by
\(\psi_{p}\) up to normalised Frobenius weight); the full
\(\psi_{p}\)-equivariance on all of \(\mathsf{A}^{(p)}\) is an
additional local constraint not automatic from (4). Both axioms are
needed and neither is redundant.

\emph{Spectral content.} The identification of the \(F\)-spectrum on
\(H^{1}\) with the nontrivial zeros of \(\xi(s)\) remains conditional
on H-P\(^{\mathrm{Ar}}\) (\Cref{v4-arith:rh:hyp:HPAr}), inherited from
\Cref{v4-arith:w8-a:thm:character-identity}. The naturality of \(F\)
established here requires only its commutation with \(H^{\bullet}(\phi)\),
not a description of its spectrum.

\section{Functoriality in \(g\)}
\label{v4-arith:w9-r5:functoriality-g}

The genus filtration
\(\mathsf{ArChAlg}_{0}\subset\mathsf{ArChAlg}_{1}\subset\cdots\) of
\Cref{v4-arith:w9-r5:def:genus-filtration} generates further
functoriality structure, tying Chapter~\ref{v4-arith:w8-a:deninger-chiral-homology}
to Chapter~\ref{v4-arith:all-genus-koszul}.

\begin{definition}[genus-\(g\) restrictions of the three functors]
\label{v4-arith:w9-r5:def:genus-restricted}
\ClaimStatusDefinitional. For every \(g\geq 0\), let
\(H^{\bullet}_{\Delta,g}\), \(H^{\bullet}_{\mathrm{ch},g}\), and
\((\partial\circ\mathrm{BdryACS})_{g}\) denote the restrictions of
\(H^{\bullet}_{\Delta}\), \(H^{\bullet}_{\mathrm{ch}}\), and
\(\partial\circ\mathrm{BdryACS}\) to \(\mathsf{ArChAlg}_{g}\).
Explicitly:
\begin{itemize}
\item \(H^{\bullet}_{\Delta,g}\) sends a genus-\(g\) object
\(\mathsf{A}\in\mathsf{ArChAlg}_{g}\) to the rank-\(2g\) Deninger
cohomology with global Frobenius \(F^{(g)}\) of
\Cref{v4-arith:w8-a:thm:all-genus-deninger}(2);
\item \(H^{\bullet}_{\mathrm{ch},g}\) sends \(\mathsf{A}\) to the
rank-\(2g\) Platonic bar homology of
\Cref{v4-arith:w8-a:def:deninger-g} with rank-\(2g\) Fock factors per
place;
\item \((\partial\circ\mathrm{BdryACS})_{g}\) sends \(\mathsf{A}\) to
the boundary chiral homology of the rank-\(2g\) 3D arithmetic CS
\(\mathrm{ACS}^{\mathrm{qu},g}_{k^{\mathrm{Ar},g}}\) of
\Cref{v4-arith:3d-acs:cor:all-genus-arith-gravity}.
\end{itemize}
\end{definition}

\begin{theorem}[compatibility under genus inclusion]
\label{v4-arith:w9-r5:thm:genus-inclusion}
\ClaimStatusProvedHereConditional \emph{(on
\Cref{v4-arith:w9-r5:thm:naturality-A},
\Cref{v4-arith:ag:thm:all-genus},
\Cref{v4-arith:w8-a:thm:all-genus-deninger})}. For every \(g\geq 0\),
the three genus-\(g\) functors commute with the inclusion
\(\iota_{g,g+1}\colon\mathsf{ArChAlg}_{g}\hookrightarrow\mathsf{ArChAlg}_{g+1}\)
up to canonical natural isomorphisms
\(\vartheta^{\Delta}_{g}\), \(\vartheta^{\mathrm{ch}}_{g}\),
\(\vartheta^{\partial}_{g}\):
\begin{equation}
\label{v4-arith:w9-r5:eq:genus-inclusion}
\begin{tikzcd}[column sep=large,row sep=large]
\mathsf{ArChAlg}_{g}\arrow[rr,"\iota_{g,g+1}"]\arrow[dr,"H^{\bullet}_{\star,g}\otimes V_{2}"']
&& \mathsf{ArChAlg}_{g+1}\arrow[dl,"H^{\bullet}_{\star,g+1}"]\\
& \mathrm{grVect}_{\mathbb{C}}^{\mathrm{Frob}}
\end{tikzcd}
\end{equation}
commutes up to \(\vartheta^{\star}_{g}\) for
\(\star\in\{\Delta,\mathrm{ch},\partial\}\), where \(V_{2}\) is the
rank-\(2\) archimedean Heisenberg summand and \(\otimes V_{2}\) is the
tensor-up with its associated graded Fock module. The three natural
transformations \(\vartheta^{\star}_{g}\) are compatible with the
natural isomorphisms \(\eta_{\Delta,\mathrm{ch}}\),
\(\eta_{\mathrm{ch},\partial}\) of
\Cref{v4-arith:w9-r5:eq:natural-iso}.
\end{theorem}

\begin{proof}
The map \(\iota_{g,g+1}\) adjoins a rank-\(2\) archimedean Heisenberg
factor at every place: at each prime \(p\), the rank-\(2g\) Fock
\(\mathcal{H}^{(p),g}\) is replaced by the rank-\(2(g+1)\) Fock
\(\mathcal{H}^{(p),g+1}=\mathcal{H}^{(p),g}\otimes\mathcal{H}^{(p)}_{2}\)
where \(\mathcal{H}^{(p)}_{2}\) is a freshly-added rank-\(2\) Fock
factor; analogously at \(\infty\).

For \(\star=\mathrm{ch}\): the Platonic bar is multiplicative under
the tensor product of factorisation algebras. At each finite stage
\(S\), the bar factors as
\(B^{\mathrm{Ar}}_{S}(\mathsf{A}\otimes V_{2})=B^{\mathrm{Ar}}_{S}(\mathsf{A})\widehat{\otimes}
B^{\mathrm{Ar}}_{S}(V_{2})\), where the right-hand factor is the bar
of the freshly-added rank-\(2\) Heisenberg. Passing to cohomology and
to the filtered colimit, the natural isomorphism
\(\vartheta^{\mathrm{ch}}_{g}\) is the K\"unneth map of the Platonic
bar, which is a quasi-isomorphism after restricting to the primary
subspace of each side (\Cref{v4-arith:thm:heisenberg-bar-zeta}
multiplicativity).

For \(\star=\Delta\): \(\vartheta^{\Delta}_{g}\) is the same map under
\Cref{v4-arith:w8-a:def:deninger-cohomology}'s definitional equality
\(H^{\bullet}_{\Delta,g}=H^{\bullet}_{\mathrm{ch},g}\); naturality is
inherited from \(\vartheta^{\mathrm{ch}}_{g}\).

For \(\star=\partial\): the bulk 3D ACS theory
\(\mathrm{ACS}^{\mathrm{qu},g+1}\) of
\Cref{v4-arith:3d-acs:cor:all-genus-arith-gravity} factors as the
tensor product of \(\mathrm{ACS}^{\mathrm{qu},g}\) with the rank-\(2\)
tensor-slot theory. This is the content of the block-diagonal splitting of the gauge
group \(GL_{2(g+1)}\supset GL_{2g}\times GL_{2}\) on the boundary:
the rank-\(2(g+1)\) WZW algebra decomposes as the tensor product of
the rank-\(2g\) WZW algebra with the rank-\(2\) WZW algebra under
this block-diagonal subgroup. The boundary chiral homology
is multiplicative under tensor products of WZW algebras, so
\(\vartheta^{\partial}_{g}\) is the K\"unneth map on the boundary
side.

Compatibility with \(\eta\). The two natural transformations
\(\vartheta^{\mathrm{ch}}_{g}\) and \(\vartheta^{\partial}_{g}\) both
come from K\"unneth maps at the corresponding bar level. Applying
the chain-level identification
\Cref{v4-arith:w9-r5:thm:chain-identification} to the rank-\(2\)
factor (whose bar is freshly introduced on both sides), we obtain
\(\eta_{\mathrm{ch},\partial,\mathsf{A}\otimes V_{2}}\circ\vartheta^{\mathrm{ch}}_{g}=\vartheta^{\partial}_{g}\circ\eta_{\mathrm{ch},\partial,\mathsf{A}}\)
as required.
\end{proof}

\begin{theorem}[uniformity: genus-wise triple identification]
\label{v4-arith:w9-r5:thm:genus-uniformity}
\ClaimStatusProvedHereConditional \emph{(on
\Cref{v4-arith:w9-r5:thm:genus-inclusion} and
\Cref{v4-arith:w8-a:thm:all-genus-deninger})}. The triple
identification is uniform in genus: for every \(g\geq 0\) and every
\(\mathsf{A}\in\mathsf{ArChAlg}_{g}\),
\begin{equation}
\label{v4-arith:w9-r5:eq:genus-uniform}
H^{\bullet}_{\Delta,g}(\mathsf{A})
\;=\;H^{\bullet}_{\mathrm{ch},g}(\mathsf{A})
\;=\;(\partial\circ\mathrm{BdryACS})_{g}(\mathsf{A})
\end{equation}
as graded vector spaces with genus-\(g\) Frobenius \(F^{(g)}\) and PD
pairing \(\mathrm{PD}^{(g)}\). The equality is compatible with the
inclusion \(\iota_{g,g+1}\) of
\Cref{v4-arith:w9-r5:thm:genus-inclusion} and with the retraction
\(\pi_{g+1,g}\): the three genus-\(g\) identifications commute, up to
\(\vartheta^{\star}_{g}\), with the corresponding genus-\((g+1)\)
identifications.
\end{theorem}

\begin{proof}
The three equalities on the right-hand side of
\eqref{v4-arith:w9-r5:eq:genus-uniform} are the genus-\(g\) restrictions
of \Cref{v4-arith:w9-r5:cor:factoring}. They hold for every
\(\mathsf{A}\in\mathsf{ArChAlg}_{g}\) by restriction.

The compatibility with genus inclusion: apply
\Cref{v4-arith:w9-r5:thm:genus-inclusion} to obtain the three
compatibility diagrams connecting \(g\) and \(g+1\). Each diagram is
commutative up to \(\vartheta^{\star}_{g}\); the three
\(\vartheta^{\star}_{g}\) agree under the natural isomorphisms
\(\eta_{\Delta,\mathrm{ch}}\) and \(\eta_{\mathrm{ch},\partial}\). Hence
the genus-\(g\) triple identification commutes with the genus-\((g+1)\)
identification: the diagram
\[
\begin{tikzcd}[column sep=scriptsize, row sep=small]
H^{\bullet}_{\Delta,g}(\mathsf{A})\otimes V_{2}\arrow[r,"="]\arrow[d,"\vartheta^{\Delta}_{g}"']
&H^{\bullet}_{\mathrm{ch},g}(\mathsf{A})\otimes V_{2}\arrow[r,"="]\arrow[d,"\vartheta^{\mathrm{ch}}_{g}"']
&(\partial\circ\mathrm{BdryACS})_{g}(\mathsf{A})\otimes V_{2}\arrow[d,"\vartheta^{\partial}_{g}"]\\
H^{\bullet}_{\Delta,g+1}(\iota(\mathsf{A}))\arrow[r,"="]
&H^{\bullet}_{\mathrm{ch},g+1}(\iota(\mathsf{A}))\arrow[r,"="]
&(\partial\circ\mathrm{BdryACS})_{g+1}(\iota(\mathsf{A}))
\end{tikzcd}
\]
commutes.

The compatibility with retraction \(\pi_{g+1,g}\) is the mirror
statement: projection onto the first \(2g\) oscillators, which is
K\"unneth-compatible by the same argument in the reverse direction.
\end{proof}

\begin{corollary}[all-genus triple functoriality]
\label{v4-arith:w9-r5:cor:all-genus-triple}
\ClaimStatusProvedHereConditional \emph{(on
\Cref{v4-arith:w9-r5:thm:genus-uniformity})}. The three functors
\(H^{\bullet}_{\Delta}\), \(H^{\bullet}_{\mathrm{ch}}\), and
\(\partial\circ\mathrm{BdryACS}\) are naturally isomorphic as
\emph{genus-filtered} functors
\(\mathsf{ArChAlg}\to\mathrm{grVect}_{\mathbb{C}}^{\mathrm{Frob}}\).
The natural isomorphism is genus-graded: restricts to a natural
isomorphism at every genus \(g\), and commutes with inclusion and
retraction along the filtration.
\end{corollary}

\begin{proof}
Assemble \Cref{v4-arith:w9-r5:cor:triple-A} (functoriality in
\(\mathsf{A}\) at fixed genus) with
\Cref{v4-arith:w9-r5:thm:genus-uniformity} (uniform compatibility
across genus). The genus-filtered structure is preserved at each
genus, with transition natural transformations \(\vartheta^{\star}_{g}\)
compatible.
\end{proof}

\subsection*{Source-disjointness and scope of genus functoriality}
\label{v4-arith:w9-r5:ah:g}

\emph{Source consistency.} The genus filtration via rank-\(2\) tensor
slots of Ch.~\ref{v4-arith:all-genus-koszul} and the 3D ACS gauge-group
scaling \(GL_{2g}\hookrightarrow GL_{2(g+1)}\) of
Ch.~\ref{v4-arith:3d-acs-E1bar} agree: in both cases the rank parameter
is \(2g\), the archimedean Heisenberg has rank \(2g\), and the Satake
parameter at each finite prime is also of rank \(2g\). The three
chapters use identical rank conventions.

\emph{Conventions.} The notation \(r=2g\) for the rank of the
archimedean Heisenberg, \(g\) for the Siegel-modular genus, and the
Satake rank \(2g\) at each finite prime is consistent across
Ch.~\ref{v4-arith:all-genus-koszul},
Ch.~\ref{v4-arith:3d-acs-E1bar}, and
Ch.~\ref{v4-arith:w8-a:deninger-chiral-homology}.

\emph{K\"unneth flatness.} The K\"unneth map for the Platonic bar at
the rank-\(2\) slot requires flatness of \(\mathcal{H}^{(p)}_{2}\)
over the Iwasawa algebra \(\Lambda_{p}\). This is the rank-\(2\)
instance of the Positselski Koszul duality of
Ch.~\ref{v4-arith:w5-d:chapter}, which holds in the abelian Fock
setting.

\emph{Retraction.} The retraction \(\pi_{g+1,g}\) does not require a
cokernel: on the primary Heisenberg scope the freshly-added rank-\(2\)
factor is a direct summand, and \(\pi_{g+1,g}\) is direct-summand
projection.

\emph{Spectral content.} The genus-\(g\) \(F\)-spectrum identification
with motivic zeros is conditional on rank-\(g\) H-P\(^{\mathrm{Ar}}\)
(\Cref{v4-arith:w8-a:thm:all-genus-deninger}(4)) and is not supplied
here. The genus-filtration structure holds independently of the
spectrum identification.

\section{Master Theorem and Explicit Example}
\label{v4-arith:w9-r5:master}

Functoriality in \(\mathsf{A}\) and genus-filtered functoriality hold
independently; the master theorem records that both hold simultaneously,
with a specific chain-level quasi-isomorphism \(\Xi\) as canonical
witness. The diagonal morphism
\(\mathsf{H}^{\mathrm{Ar}}\to\mathsf{H}^{\mathrm{Ar}}\oplus\mathsf{H}^{\mathrm{Ar}}\)
exercises every axiom and every piece of the triple identification.

\begin{theorem}[master functoriality]
\label{v4-arith:w9-r5:thm:master-functoriality}
\ClaimStatusProvedHereConditional \emph{(on the residuals of
Chapters~\ref{v4-arith:3d-acs-E1bar},
\ref{v4-arith:all-genus-koszul},
\ref{v4-arith:w8-a:deninger-chiral-homology} and their arithmetic
Positselski, H-P\(^{\mathrm{Ar}}\), 3D quantum ACS scopes)}. The
triple identification
\(H^{\bullet}_{\Delta}=H^{\bullet}_{\mathrm{ch}}=\partial\circ\mathrm{BdryACS}\)
is a natural isomorphism of genus-filtered functors
\[
\mathsf{ArChAlg}=\bigcup_{g\geq 0}\mathsf{ArChAlg}_{g}
\longrightarrow\mathrm{grVect}_{\mathbb{C}}^{\mathrm{Frob}},
\]
which is:
\begin{enumerate}
\item functorial in \(\mathsf{A}\): every morphism
\(\phi\colon\mathsf{A}_{1}\to\mathsf{A}_{2}\) in \(\mathsf{ArChAlg}\)
induces a commuting triple-square of cohomology maps, preserving
Frobenius \(F\) and Poincar\'e duality \(\mathrm{PD}\);
\item functorial in \(g\): for every \(g\geq 0\), the restriction to
\(\mathsf{ArChAlg}_{g}\) is a natural isomorphism of genus-\(g\)
functors, compatible with the inclusion
\(\iota_{g,g+1}\) and the retraction \(\pi_{g+1,g}\);
\item chain-level canonical: the natural isomorphism is induced by
the quasi-isomorphism \(\Xi\) of
\Cref{v4-arith:w9-r5:thm:chain-identification}, a
specific chain map built from arithmetic Positselski bar-cobar,
the \(E_{1}\)-chiral archimedean bar, and the Parshin--Gersten tame-symbol
layer.
\end{enumerate}
\end{theorem}

\begin{proof}
Part~(1) is \Cref{v4-arith:w9-r5:cor:triple-A}.
Part~(2) is \Cref{v4-arith:w9-r5:cor:all-genus-triple}.
Part~(3) is the explicit construction of \(\Xi\) in the proof of
\Cref{v4-arith:w9-r5:thm:chain-identification}, together with the
fact that the Deninger-to-chiral-homology equality
\(\eta_{\Delta,\mathrm{ch}}\) is definitional
(\Cref{v4-arith:w8-a:def:deninger-cohomology}) and hence trivially
chain-level.
\end{proof}

\begin{example}[diagonal \(\mathsf{H}^{\mathrm{Ar}}\to\mathsf{H}^{\mathrm{Ar}}\oplus\mathsf{H}^{\mathrm{Ar}}\)]
\label{v4-arith:w9-r5:ex:diagonal}
The diagonal morphism
\[
\Delta\colon\mathsf{H}^{\mathrm{Ar}}\longrightarrow
\mathsf{H}^{\mathrm{Ar}}\oplus\mathsf{H}^{\mathrm{Ar}}
\]
is defined place-wise by the arithmetic Heisenberg diagonal
\[
\Delta^{(v)}\colon\mathcal{H}^{(v)}\longrightarrow\mathcal{H}^{(v)}\oplus\mathcal{H}^{(v)},
\qquad x\longmapsto(x,x),
\]
at every place \(v\in|\overline{C}|\). The target is the coproduct of
two copies of \(\mathsf{H}^{\mathrm{Ar}}\) in \(\mathsf{ArChAlg}\); its
archimedean component is the free \(E_{1}\)-chiral product of two
rank-\(2\) Heisenbergs, hence a rank-\(4\) Heisenberg.
\end{example}

\begin{proposition}[triple verification on the diagonal]
\label{v4-arith:w9-r5:prop:diagonal-verification}
\ClaimStatusProvedHereConditional \emph{(on
\Cref{v4-arith:w9-r5:thm:master-functoriality})}. The diagonal
\(\Delta\) of \Cref{v4-arith:w9-r5:ex:diagonal} is a morphism in
\(\mathsf{ArChAlg}\) (axioms (1)--(6) of
\Cref{v4-arith:w9-r5:def:ArChAlg-morphism} are satisfied), and
the induced triple-morphism
\[
\bigl(H^{\bullet}_{\Delta}(\Delta),\,H^{\bullet}_{\mathrm{ch}}(\Delta),\,\partial\circ\mathrm{BdryACS}(\Delta)\bigr)
\]
commutes with \(\eta_{\Delta,\mathrm{ch}}\) and
\(\eta_{\mathrm{ch},\partial}\). Explicitly:
\begin{enumerate}
\item \(H^{\bullet}_{\Delta}(\Delta)\) is the diagonal embedding
\(H^{\bullet}_{\Delta}(\mathsf{H}^{\mathrm{Ar}})\to H^{\bullet}_{\Delta}(\mathsf{H}^{\mathrm{Ar}}\oplus\mathsf{H}^{\mathrm{Ar}})=H^{\bullet}_{\Delta}(\mathsf{H}^{\mathrm{Ar}})\oplus H^{\bullet}_{\Delta}(\mathsf{H}^{\mathrm{Ar}})\)
of the arithmetic cohomology; in degree \(0\) it is \(x\mapsto(x,x)\)
on \(\mathbb{C}_{0}\to\mathbb{C}_{0}\oplus\mathbb{C}_{0}\); similarly
for \(H^{2}\); and it sends the middle \(H^{1}\)-eigenvector of the Weil-operator
\(\Theta_{\xi}\) (the functional-equation involution of
\Cref{v4-arith:rh:def:ThetaXi}) at ordinate \(\gamma_{\rho}\) to the
pair of two copies of the same eigenvector.
\item \(H^{\bullet}_{\mathrm{ch}}(\Delta)\) is the induced chain map
of Platonic bars, compatible with the bar differential because
\(\Delta\) preserves the spherical units and is a homomorphism of
chiral algebras.
\item \(\partial\circ\mathrm{BdryACS}(\Delta)\) is the boundary
restriction of the bulk interface
\(\mathrm{BdryACS}(\Delta)\colon\mathrm{ACS}^{\mathrm{qu}}_{\mathsf{H}^{\mathrm{Ar}}}\to\mathrm{ACS}^{\mathrm{qu}}_{\mathsf{H}^{\mathrm{Ar}}}\oplus\mathrm{ACS}^{\mathrm{qu}}_{\mathsf{H}^{\mathrm{Ar}}}\),
which in the 3D ACS framework is the diagonal bulk interface defect connecting one bulk theory to the
disjoint union of two copies of itself.
\end{enumerate}
All three commute with the natural isomorphisms
\(\eta_{\Delta,\mathrm{ch}}\), \(\eta_{\mathrm{ch},\partial}\).
\end{proposition}

\begin{proof}
Axiom verification. Axiom (1): \(\Delta^{(\infty)}\) is the diagonal
on the archimedean Heisenberg factorisation algebra; it preserves
the conformal vector because the free product of chiral algebras has
a conformal vector which is the sum of the two factor conformal
vectors, and the diagonal \(T\mapsto T\oplus T\) is consistent with
the coproduct structure. Axiom (2): \(\Delta^{(p)}\) commutes with
\(\psi_{p}\) because \(\psi_{p}(x,x)=(\psi_{p}(x),\psi_{p}(x))\). Axiom
(3): \(\Delta^{(v)}(u^{(v)})=(u^{(v)},u^{(v)})\), which is the
spherical unit of the coproduct. Axiom (4): the adjacent product
on the coproduct is the diagonal of the adjacent product on each
summand; \(\Delta\) commutes with this by construction. Axiom (5): the
tame-symbol boundary \(\partial^{\mathrm{PG}}\) is linear, so
\(\partial^{\mathrm{PG}}(x,x)=(\partial^{\mathrm{PG}}(x),\partial^{\mathrm{PG}}(x))\).
Axiom (6): Tate compatibility is preserved because the spherical
units are preserved and the restricted tensor factors through each
summand.

Claim (1). The Deninger cohomology of the coproduct
\(\mathsf{H}^{\mathrm{Ar}}\oplus\mathsf{H}^{\mathrm{Ar}}\) is
\(H^{\bullet}_{\Delta}(\mathsf{H}^{\mathrm{Ar}})\oplus H^{\bullet}_{\Delta}(\mathsf{H}^{\mathrm{Ar}})\)
because arithmetic bar commutes with direct sums on the primary
scope (the bar is additive in its argument, a consequence of
\Cref{v4-arith:platonic:def:bar} and of the additivity of the
tame-symbol boundary). The induced map is the diagonal. In each
degree:
\(H^{0}=\mathbb{C}_{0}\to\mathbb{C}_{0}\oplus\mathbb{C}_{0}\),
\(x\mapsto(x,x)\); \(H^{2}=\mathbb{C}_{1}\to\mathbb{C}_{1}\oplus\mathbb{C}_{1}\),
\(x\mapsto(x,x)\); \(H^{1}\)-eigenvectors at ordinate \(\gamma_{\rho}\)
go to the pair of two copies.

Claim (2). The Platonic bar is natural in \(\mathsf{A}\) by
\Cref{v4-arith:w9-r5:def:functor-ch}; applying naturality to
\(\Delta\) gives the chain map on bars. Bar differential compatibility
is axiom (5) of
\Cref{v4-arith:w9-r5:def:ArChAlg-morphism} applied to the diagonal.

Claim (3). The bulk-boundary assignment
\(\mathrm{BdryACS}(\Delta)\) is by
\Cref{v4-arith:w9-r5:def:functor-bdryACS} the 3D bulk interface
whose archimedean boundary restriction is \(\Delta^{(\infty)}\) (the
diagonal map of archimedean \(E_{1}\)-chiral algebras) and whose
finite-place Wilson-line restriction is
\(\{\Delta^{(p)}\}\). The boundary chiral homology \(\partial\) of this
interface is the induced map on boundary chiral homology. By
the additivity of boundary chiral homology for coproducts in
\(\mathsf{3DACS}_{\overline{\mathrm{Spec}\,\mathbb{Z}}}\) (the disjoint
union of two bulk theories has as boundary the coproduct of
the two boundary chiral algebras), the boundary map is the diagonal
on the rank-\(2\) chiral homology.

Compatibility with \(\eta\). By construction each of the three maps
agrees with the others under the natural isomorphisms: they are all
the diagonal of the same underlying primary Heisenberg cohomology
object, related by the chain-level quasi-isomorphism \(\Xi\) of
\Cref{v4-arith:w9-r5:thm:chain-identification} on each direct
summand and the identity on the direct-sum structure.
\end{proof}

\begin{remark}[what the diagonal example witnesses]
\label{v4-arith:w9-r5:rem:diagonal-witnesses}
The diagonal \(\Delta\) is the simplest non-trivial morphism in
\(\mathsf{ArChAlg}\) after the identity: it is not an isomorphism,
not a projection, not a single-factor inclusion. Verifying the
triple-morphism commutation on \(\Delta\) tests every axiom of
\Cref{v4-arith:w9-r5:def:ArChAlg-morphism} against every piece of
the triple identification. Because the coproduct
\(\mathsf{H}^{\mathrm{Ar}}\oplus\mathsf{H}^{\mathrm{Ar}}\) sits in
\(\mathsf{ArChAlg}_{g=2}\) (rank-\(4\) archimedean Heisenberg) while
\(\mathsf{H}^{\mathrm{Ar}}\) sits in \(\mathsf{ArChAlg}_{g=1}\), the
example also touches genus functoriality: it is a non-trivial map
in the filtered structure
\(\mathsf{ArChAlg}_{1}\to\mathsf{ArChAlg}_{2}\) that does not factor
through the filtered-colimit inclusion \(\iota_{1,2}\).
\end{remark}

\subsection*{Source-disjointness and scope of the master theorem}
\label{v4-arith:w9-r5:ah:master}

\emph{Source-disjointness.} The master theorem combines three source
families: Beilinson--Drinfeld 2004 (archimedean \(E_{1}\)-chiral
structure), Fargues--Scholze arXiv:2102.13459 (finite-prime
perfectoid charts), and Kim 2015 arXiv:1510.05818 plus the quantum
ACS (3D bulk theory). These families are pairwise disjoint in the
Vol~IV HZ-IV sense.

\emph{Sign consistency.} The Koszul signs in the Platonic bar
differential, the alternating signs in the ordered-incidence bar,
and the 3D framed-cobordism signs in the boundary-chiral-homology
pairing are mutually consistent under the chain-level identification
\(\Xi\) of \Cref{v4-arith:w9-r5:thm:chain-identification}.

\emph{Inherited conditional hypotheses.} The well-definedness of
\(\partial\circ\mathrm{BdryACS}\) as a functor (rather than a mere
correspondence) requires a well-defined bulk interface defect
category, conditional on
\Cref{v4-arith:3d-acs:thm:finite-bulk-boundary} and the renormalised
profinite limit of \Cref{v4-arith:3d-acs:conj:qu-acs}. Both
conditionals are inherited from Chapter~\ref{v4-arith:3d-acs-E1bar}.

\emph{Spectral scope.} The master theorem claims the
\emph{functoriality} of the triple identification, not the spectral
content. The three spectral statements (Deninger purity, Riemann
Hypothesis, \(F\)-RH) remain conditional on H-P\(^{\mathrm{Ar}}\) and
the archimedean A-TP hypothesis; neither is asserted here.

\section{Scope of the Master Theorem}
\label{v4-arith:w9-r5:residual-summary}

\subsection{Functoriality residuals of Chapter~\ref{v4-arith:w8-a:deninger-chiral-homology} closed here}
\label{v4-arith:w9-r5:residual-closed}

The following residuals of
Chapter~\ref{v4-arith:w8-a:deninger-chiral-homology} are closed here:
\begin{enumerate}
\item \Cref{v4-arith:w8-a:residual}(5): triple-identification
functoriality under pullback along number-field extensions
\(K/\mathbb{Q}\). Closed: any such pullback is a morphism in
\(\mathsf{ArChAlg}\) (it preserves spherical units, adjacent products,
and Parshin--Gersten boundaries by classical number-field theory),
and functoriality follows from
\Cref{v4-arith:w9-r5:cor:triple-A}.
\item \Cref{v4-arith:w8-a:residual}(6): genus-\(g\) Wilson-line
Frobenius identification. Closed: the genus-\(g\) restriction of
\(\partial\circ\mathrm{BdryACS}\) is functorially identified with
the genus-\(g\) Deninger cohomology by
\Cref{v4-arith:w9-r5:thm:genus-uniformity}; the rank-\(2g\)
Wilson-loop observables are identified with the rank-\(2g\)
\(\psi_{p}^{(g)}\) by naturality in the direct summand structure.
\end{enumerate}

\subsection{Inherited conditional hypotheses of the master theorem}
\label{v4-arith:w9-r5:residual-inherited}

The master functoriality theorem
\Cref{v4-arith:w9-r5:thm:master-functoriality} is conditional on
the following inherited hypotheses:
\begin{enumerate}
\item Arithmetic Positselski (Ch.~\ref{v4-arith:w5-d:chapter}) for
the finite-prime factor quasi-isomorphism \(\Xi^{(p)}\).
\item 3D quantum ACS existence
(\Cref{v4-arith:3d-acs:conj:qu-acs}) for the well-definedness of
\(\mathrm{BdryACS}\) as a functor.
\item Local-trace comparison
(\Cref{v4-arith:3d-acs:thm:bulk-boundary-Vprim}) for the boundary
identification at \(\infty\).
\item H-P\(^{\mathrm{Ar}}\)
(\Cref{v4-arith:rh:hyp:HPAr}) for the spectrum identification on
\(H^{1}\); not used in functoriality but used in spectral content.
\item Rank-\(g\) K\"unneth compatibility on the Platonic bar, which
is standard by rank-\(g\) Positselski
(Ch.~\ref{v4-arith:ag:residual}(2)).
\end{enumerate}
No new residuals are introduced.

\subsection{The master theorem and what it forces}
\label{v4-arith:w9-r5:summary}

The triple identification
\(H^{\bullet}_{\Delta}=H^{\bullet}_{\mathrm{ch}}=\partial\circ\mathrm{BdryACS}\),
previously established object-wise for \(\mathsf{H}^{\mathrm{Ar}}\)
at genus \(g=0\), is a natural isomorphism of genus-filtered functors
on all of \(\mathsf{ArChAlg}\): every morphism of arithmetic chiral
algebras induces a commuting triple-square of cohomology maps,
Frobenius-equivariant and Poincar\'e-dual-compatible; every
genus-inclusion and retraction is intertwined by K\"unneth
natural transformations \(\vartheta^{\star}_{g}\). The chain-level
quasi-isomorphism \(\Xi\) of
\Cref{v4-arith:w9-r5:thm:chain-identification}, built from arithmetic
Positselski bar-cobar at finite primes, the \(E_{1}\)-chiral archimedean
bar, and the Parshin--Gersten tame-symbol layer, is the canonical
chain-level witness. The three source families (Deninger 1998,
Beilinson--Drinfeld + Positselski, Kim + quantum ACS) are pairwise
disjoint; the natural isomorphism is therefore a genuine
Vol~IV HZ-IV verification. The remaining spectral content (that the Frobenius \(F\) on \(H^{1}\)
has eigenvalues at the nontrivial zeros of \(\xi(s)\)) requires the
inherited H-P\(^{\mathrm{Ar}}\) hypothesis
(\Cref{v4-arith:rh:hyp:HPAr}) and is not supplied here.

\section*{Primary Sources and Disjointness}
\label{v4-arith:w9-r5:bibliography}

The chapter depends on four disjoint source families.

\begin{description}
\item[Arithmetic chiral algebra side.]
Beilinson--Drinfeld 2004, \emph{Chiral Algebras} (AMS Colloq.~51);
Francis 2013 (geometric \(E_{1}\)-ordered bar);
Francis--Gaitsgory 2012 arXiv:1103.5803 (factorisation gluing);
Costello--Gwilliam 2017, \emph{Factorization Algebras in Quantum
Field Theory} (Cambridge UP) for the archimedean \(E_{1}\)-chiral
factorisation algebra.
\item[Finite-prime arithmetic side.]
Fargues--Scholze arXiv:2102.13459 (geometrisation of local Langlands);
Bhatt--Scholze 2022, \emph{Prisms and prismatic cohomology}, Ann.~Math.
for the prismatic/Nygaard comparison; Borger 2009 arXiv:0906.3146
for the Bost--Connes Adams \(\psi_{p}\); Positselski 2011, Mem.~AMS~996,
\emph{Two kinds of derived categories, Koszul duality, and
comodule-contramodule correspondence} for the finite-prime
quasi-isomorphism \(\Xi^{(p)}\).
\item[3D ACS / bulk-boundary side.]
Kim 2015 arXiv:1510.05818, \emph{Arithmetic Chern--Simons Theory~I};
Chern--Dan--Kim--Park 2018 arXiv:1609.03012 (non-abelian extension);
Mazur 1973 ASENS~6 (\'etale cohomological dimension of number rings);
Morishita 2012 arXiv:0904.3399, \emph{Knots and Primes};
Witten 1989, CMP~121, and Witten 2007 arXiv:0706.3359 (3D gravity /
3D CS boundary WZW); Artin--Verdier duality per SGA 4.5.
\item[Genus and modular side.]
Consani--Marcolli 2006 (genus-\(g\)
Bost--Connes--Marcolli); Faltings 1984, Ann.~Math.~119, \emph{Calculus
on arithmetic surfaces}; Kudla 1997, Ann.~Math.~146; Li--Liu 2021
arXiv:2101.09485 (Kudla--Rapoport at bad primes); Arthur 2013,
\emph{The Endoscopic Classification of Representations}, AMS Colloq.~61.
\end{description}

Pairwise disjointness of the four families is witnessed by the
Vol~IV source-catalog
\Cref{v4-arith:app:A:langlands}--\Cref{v4-arith:app:A:costello},
\Cref{v4-arith:app:A:prismatic},
\Cref{v4-arith:app:A:arith-Ddagger},
and the 3D ACS citation cluster of
Ch.~\ref{v4-arith:3d-acs-E1bar}. Each theorem of this chapter is
decorable in the Vol~IV HZ-IV sense: derived path through one
family, verified against a second family.
